\section{Індивідуальне завдання}

В межах індивідуального завдання за варіантом 1 досліджено асиметричний криптографічний механізм
ElGamal.

\subsection{Призначення та задача}

Криптосистема ElGamal є асиметричним криптографічним алгоритмом, який було розроблено Тахером Ель
Гамалем у 1985 році. Вона призначена для розв'язання двох основних задач захисту інформації: забезпечення
конфіденційності даних шляхом асиметричного шифрування та забезпечення автентичності й цілісності
повідомлень за допомогою цифрового підпису.

\subsection{Основна ідея роботи}

Основна ідея криптосистеми ElGamal базується на використанні односторонньої функції з таємним входом у
мультиплікативній групі скінченного поля. Безпека системи повністю спирається на обчислювальну складність
задачі дискретного логарифмування. Під час шифрування відправник використовує відкритий ключ
одержувача та генерує одноразовий випадковий сеансовий параметр. Це забезпечує властивість
ймовірнісного шифрування, коли один і той самий відкритий текст при кожному повторному шифруванні дає
новий шифротекст, що унеможливлює атаки на основі аналізу повторюваності даних.

\subsection{Основні криптографічні примітиви}

В основі функціонування криптосистеми ElGamal лежать операції в мультиплікативних групах скінченних
полів за модулем великого простого числа. Алгоритм спирається на складність обчислення дискретного
логарифма у скінченному полі, застосовує односторонні криптографічні хеш-функції для формування
цифрового підпису та вимагає використання криптографічно стійких генераторів псевдовипадкових чисел
для вибору одноразового сеансового параметра.

\subsection{Властивості безпеки}

Криптосистема ElGamal забезпечує конфіденційність інформації, оскільки розшифрувати повідомлення може
лише власник відповідного приватного ключа. Завдяки ймовірнісному шифруванню алгоритм гарантує
непередбачуваність шифротексту. У режимі цифрового підпису система забезпечує автентичність джерела та
невідмовність від авторства, а також контролює цілісність даних, унеможливлюючи їх непомітну модифікацію
під час передачі.

\subsection{Приклад практичного сценарію застосування}

Математичні принципи ElGamal лежать в основі стандарту цифрового підпису DSA, а також застосовуються у
протоколах захисту електронної пошти OpenPGP та GnuPG для асиметричного шифрування сеансових ключів
і підписання повідомлень. Крім того, гомоморфні властивості шифру ElGamal дозволяють виконувати
обчислення над зашифрованими даними без їх розшифрування, що широко використовується у протоколах
електронного голосування для підрахунку голосів.

\subsection{Основні обмеження та ризики практичного використання}

Основним недоліком ElGamal є надмірність шифротексту, розмір якого удвічі перевищує розмір початкового
відкритого тексту. Повторне використання одного й того самого випадкового сеансового параметра для
шифрування різних повідомлень призводить до розкриття відкритого тексту. Без використання спеціальних
схем доповнення ElGamal є податливим до атак на основі мультиплікативної гомоморфності, а самі модульні
обчислення вимагають значно більше ресурсів порівняно із сучасними алгоритмами на еліптичних кривих.
